\label{annexe:orchestrator}

Extrait du cœur logique de l'agent, montrant la gestion de la mémoire et la stratégie de repli (Fallback) entre le Graphe et le Vecteur.

\begin{lstlisting}[language=python, caption=Algorithme de décision (orchestrator.py - Extrait)]
async def process_user_message(text: str, session_id: str):
    
    # 1. Gestion de la Memoire (Redis)
    memory = ConversationMemory(session_id)
    history = memory.get_history()
    
    # Reformulation contextuelle via Gemini
    standalone_text = text
    if history:
        standalone_text = contextualizer.rewrite(history, text)

    # 2. Strategie GraphRAG (Prioritaire)
    print("Tentative 1 : GraphRAG (Neo4j)...")
    graph_answer = graph_engine.query(standalone_text)
    
    if graph_answer and "je ne sais pas" not in graph_answer.lower():
        final_response = graph_answer
    
    # 3. Strategie VectorRAG (Fallback)
    else:
        print("GraphRAG muet. Tentative 2 : VectorRAG...")
        rag_results = rag_engine.search(standalone_text)
        
        if rag_results and rag_results[0]['score'] > 0.25:
            final_response = rag_results[0]['content']
        else:
            final_response = "Je ne trouve pas l'information."

    # 4. Sauvegarde
    memory.add_message("user", text)
    memory.add_message("ai", final_response)
    
    return final_response
\end{lstlisting}
